\section{Construction}\label{sec:construction}

We first build intuition for the construction by describing, informally, a
sequence of increasingly refined solutions, before giving the protocol in
full.

\subsection{Design overview}

In its simplest form, a token is a triple $\tau = (\pk, v, s)$, where $\pk$
is the public key of its owner, $v$ is its value, and $s$ is a random number
that uniquely identifies it.

\paragraph{Withdrawals.} User $\usr_0$ sends her intermediary $I_0$ a
withdrawal request $\langle v, s_0 \rangle$, where $v$ is the value to be
withdrawn and $s_0$ is a random number she selects. After checking that
$\usr_0$ has enough funds, $I_0$ sends the triple $\langle \pk_0, v, s_0
\rangle$ to the central bank, which checks that $I_0$ has enough funds and
responds with signature $\psi_0$ on $\langle \pk_0, v, s_0 \rangle$. With
$\psi_0$, $\usr_0$ can spend token $\tau_0 = (\pk_0, v, s_0)$. A payment in
this scheme transfers ownership of a token from a payer to a payee; it never
splits or merges tokens.

\paragraph{Payments.} To pay user $\usr_1$ with $\tau_0$, $\usr_0$ computes a
signature $\sigma_0$ on $(v, s_0, \pk_0, \pk_1)$ and sends the tuple
$\langle v, s_0, \pk_0, \pk_1, \sigma_0, \psi_0 \rangle$. $\usr_1$ accepts
the payment if $\sigma_0$ and $\psi_0$ both verify against this information,
and can subsequently transfer the token to a user $\usr_2$ in the same way.
After $n$ consecutive payments, the $n$-th payee $\usr_n$ receives $\langle
v, s_0, \pk_0, \ldots, \pk_n, \sigma_0, \ldots, \sigma_{n-1}, \psi_0
\rangle$, and accepts it if every $\sigma_i$ is a valid signature on $(v,
s_0, \pk_i, \pk_{i+1})$ under $\pk_i$, and $\psi_0$ is a valid central bank
signature on $(\pk_0, v, s_0)$.

\paragraph{Deposits.} To deposit the token, $\usr_n$ computes a signature
$\sigma_n$ on $(v, s_0, \pk_n)$ and sends $\langle v, s_0, \pk_0, \ldots,
\pk_n, \sigma_0, \ldots, \sigma_n, \psi_0 \rangle$ to her intermediary
$I_n$, which forwards the deposit to the central bank. The central bank
verifies that $\sigma_n$ is a valid signature on $(v, s_0, \pk_n)$, that
$\psi_0$ is a valid signature of the central bank on $(\pk_0, v, s_0)$, and
that every $\sigma_i$ is a valid signature on $(v, s_0, \pk_i, \pk_{i+1})$
under $\pk_i$. If all checks pass, the deposit is accepted; if another
deposit request already starts with $(v, s_0, \pk_0)$, the central bank
rejects it as double spending.

\paragraph{Identifying double spenders.} This simple scheme is prone to
framing: naively, the double spender is identified as the public key
$\pk_j$ at the point where two conflicting deposit chains diverge, but this
assumes a double spender never spends the same token twice to the same
payee. A stateful payee who remembers every payment they have been involved
in could reject duplicates, but this is unreasonable to expect. Instead we
have each payee select a random number that uniquely identifies the payment
they are about to receive and that is carried forward in subsequent
payments, following an approach also used in transferable
e-cash~\cite{Camenisch2005,Fuchsbauer2009,Baldimtsi2015}. This ensures that
two payments from the same payer to the same payee using the same token
still carry distinct random numbers, so a payee who received the same token
twice is not at risk of being framed. Concretely, $\usr_i$ paying $\usr_{i+1}$
receives a fresh random number $s_{i+1}$ from $\usr_{i+1}$, computes
$\sigma_i$ on $(v, s_i, s_{i+1}, \pk_i, \pk_{i+1})$, and sends $\langle v,
s_0, \pk_0, s_1, \pk_1, \ldots, s_{i+1}, \pk_{i+1}, \sigma_0, \ldots,
\sigma_i, \psi_0 \rangle$ to $\usr_{i+1}$; at deposit, $\usr_n$ computes
$\sigma_n$ on $(v, s_n, \pk_n)$ and sends the analogous deposit request.
Given two conflicting deposit requests both spending $(v, s_0, \pk_0)$, the
double spender is now the user with public key $\pk_j$ such that $(\pk_i,
s_i) = (\pk'_i, s'_i)$ for all $i \leq j$ but $(\pk_{j+1}, s_{j+1}) \neq
(\pk'_{j+1}, s'_{j+1})$: this user provably sent two distinct payments, one
carrying $s_{j+1}$ and the other $s'_{j+1}$.

\paragraph{Adding anonymity.} The scheme so far achieves token
unforgeability, double spending detection, and double spender
identification, but not privacy: the chain of payments reveals every payer's
and payee's public key. We add anonymity by having users obtain long-term
credentials from the registration authority that bind their unique
identifier to their secret key, for instance using BBS+
signatures~\cite{Camenisch2016}. Rather than a public key, a token's owner is
now identified by a pseudonym: a hiding commitment to the owner's unique
identifier, so a token is $\tau = (v, C)$ with $C = \gen_0^{\uid}
\gen_1^{r}$ for random $r$. During payments and deposits, the user produces
an anonymous signature by proving in zero knowledge that the committed
identifier was signed by the registration authority, together with a
signature of knowledge that they know the corresponding secret key. This
gives a chain of payments of the form $\langle v, s_0, C_0, \ldots, s_n,
C_n, \sigma_0, \ldots, \sigma_n, \psi_0 \rangle$, where $\psi_0$ now signs
$\tau_0 = (C_0, v, s_0)$ rather than $(\pk_0, v, s_0)$. We can further
shrink this chain by not revealing $s_i$ and instead using it directly as
the randomness of $C_i = \gen_0^{\uid_i}\gen_1^{s_i}$, giving payments of the
form $\langle v, C_0, \ldots, C_n, \sigma_0, \ldots, \sigma_n, \psi_0
\rangle$.

\paragraph{Adding auditability.} To identify double spenders at deposit
time, each payment carries a \emph{verifiable} ElGamal encryption of the
payer's unique identifier, intended for an authorised auditor. Given two
deposit requests $\langle v, C_0, \cipher_0, \ldots, C_n, \cipher_n,
\sigma_0, \ldots, \sigma_n, \psi_0 \rangle$ and $\langle v, C_0, \cipher'_0,
\ldots, C'_n, \cipher'_n, \sigma'_0, \ldots, \sigma'_n, \psi_0 \rangle$
starting from the same $(v, C_0)$, the double spender is identified by
decrypting the ciphertext $\cipher_j$ at the index $j$ where $C_i = C'_i$
for all $i \leq j$ but $C_{j+1} \neq C'_{j+1}$. Since the auditor can
de-anonymise payments, we recommend distributing it via \textbf{threshold
decryption}, so decryption succeeds only if more than a threshold $t$ of
auditors cooperate.

\paragraph{Hiding the link between withdrawals and deposits.} The scheme so
far does not prevent linking a withdrawal to a deposit. This is a genuine
privacy concern even though payments themselves are anonymised, since users
identify themselves to their intermediaries at both withdrawal and deposit:
if a token is transferred only once before being deposited, and the payer's
and payee's intermediaries collude, they can de-anonymise the payment. We
avoid this using Pointcheval-Sanders blind
signatures~\cite{Pointcheval2018,Sonnino2019}, which let the central bank
sign the token $\tau_0 = (v, C_0)$ without learning $C_0$ or the resulting
signature $\psi_0$. We note that the withdrawal and deposit of a token can
still be linked through its value; this can be avoided by using a single
denomination, or mitigated by using fixed denominations together with
aggregated deposits, which let a depositor disclose only the aggregate value
of several deposited tokens rather than the value of each. We take the
latter approach (cf.\ Appendix~\ref{sec:aggregate-deposits}).

\paragraph{Using secure elements to limit double spending.} As discussed
in the introduction, secure elements make double spending attacks
impractical to mount. Since today's secure elements have limited compute and
storage, we split the computation above into a lightweight part executed by
the secure element and a heavier part offloaded to the user's mobile device.
The secure element is responsible for three tasks: generating the randomness
$s_i$ and computing the commitment $C_i$, producing the signature of
knowledge of the user's secret key, which only it can access, and deleting a
token once it has been spent. The mobile device handles everything else, in
particular verifying received payments, computing the verifiable encryption,
and generating the zero-knowledge proofs that show the user is enrolled in
the system.

\paragraph{Privacy guarantees.} To summarise: (1) transactors are anonymous
to anyone who sees a transaction transcript, assuming payer and payee
exchange identities out of band; (2) withdrawals and deposits are unlinkable
to every party in the system, so no one can trace a deposit or payment to
its origin or destination from the payment transcripts; and (3) the
auditor may deanonymise payments, and is trusted to do so only with evidence
of double spending, with the auditor role recommended to be distributed via
threshold decryption.

\subsection{The protocol in full}

We now describe the complete protocol, using the notation $\Pi \leftarrow
\pok[\instance, \wit : \rel(\instance, \wit) = 1]$ for a non-interactive
zero-knowledge proof of knowledge, and $\sigma \leftarrow \sok[m, (\instance,
\wit) : \rel(\instance, \wit) = 1]$ for a signature of knowledge on message
$m$.

\paragraph{Setup.}\label{sec:offline:construction:setup} The central bank
generates keys for a blind signature scheme supporting three-message
vectors, obtaining $(\cbsk, \cbpk)$, and runs $\ck \leftarrow
\ckeygen(1^\kappa)$ to obtain a public commitment key. The auditor generates
an asymmetric encryption key pair $(\audsk, \audpk) \leftarrow
\ekeygen(1^\kappa)$. The registration authority generates keys for a
multi-message signature scheme supporting two-message vectors, obtaining
$(\rask, \rapk) \leftarrow \skeygen(1^\kappa, 2)$. All public keys and
generators are published to every participant.

When a user $\usr$ opens an account, their secure element $\se$ is
provisioned with a fresh secret key $\usk$ and a unique random identifier
$\uid$. The registration authority verifies that $\uid$ is fresh and that
$\se$ knows $\usk$ via an interactive zero-knowledge proof of knowledge,
then issues credential $\cred = (\uid, \upk, \varphi)$ with $\varphi
\leftarrow \sigsign(\rask, (\uid, \usk))$. $\dev$ and $\se$ are paired so
that only $\dev$ can send instructions to $\se$.

\paragraph{Withdrawal.}\label{sec:withdrawal} User $\usr_0$ withdraws a token
of value $v_0$ as follows.

\emph{Preparation}\label{sec:stepW1W2}: $\dev_0$ instructs $\se_0$ to
initiate a withdrawal of value $v_0$. $\se_0$ selects random $s_0$, computes
$\comm_0 \leftarrow \com(\ck, \uid_0; s_0)$, and returns $(s_0, \comm_0)$ to
$\dev_0$, which computes $\comm \leftarrow \com(\ck, (\uid_0, v_0); s_0)$.
$\dev_0$ prepares a blind signature request on the committed vector
$(\uid_0, v_0, s_0)$, hiding $\uid_0$ and $s_0$ and disclosing $v_0$,
including a zero-knowledge proof $\Pi$ that $\comm$ is correctly formed and
that the accompanying ciphertexts are correct, and sends $\req = (\uid_0,
\comm, \cipher_0, v_0, \cipher_2, \Pi, \rho_0)$ to $\Int_0$.

\emph{Issuance}\label{sec:stepW3}: $\Int_0$ verifies $\Pi$, locks $v_0$ in
$\usr_0$'s account, and forwards $(\comm, \cipher_0, v_0, \cipher_2, \Pi)$ to
the central bank, which checks that $\comm$ is fresh, verifies $\Pi$, debits
$\Int_0$'s reserves, and returns a blind signature response $\resp$.

\emph{Receipt}\label{sec:stepW4}: $\dev_0$ unblinds $\resp$ to obtain a
signature $\psi \leftarrow \sigsign(\cbsk, (\uid_0, v_0); s_0)$ and stores
$(v_0, s_0, \comm_0, \psi)$; $\se_0$ stores the original token $\tau_0 =
(v_0, \comm_0)$.

\emph{Token creation}\label{sec:stepW5}: $\dev_0$ produces a proof of
withdrawal $\Sigma$ with $\instance = (\ck, \cbpk, v_0, \comm_0)$, $\wit =
(\uid_0, s_0, \psi)$, and relation
\begin{equation*}
  \sigverify(\cbpk, (\uid_0, v_0); s_0, \psi) = 1 \;\wedge\; \opn(\ck,
  \uid_0; s_0, \comm_0) = 1,
\end{equation*}
and stores $T_0 = (v_0, s_0, \comm_0, \hist_0 = \Sigma)$.

\paragraph{Payment.}\label{sec:payments} $\usr_i$ pays $\usr_{i+1}$ using
token $(v, s_i, \comm_i, \hist_i)$ as follows.

\emph{Preparation}\label{sec:stepS1}: $\dev_{i+1}$ instructs $\se_{i+1}$ to
prepare to receive a payment of value $v$. $\se_{i+1}$ selects random
$s_{i+1}$, computes $\comm_{i+1} \leftarrow \com(\ck, \uid_{i+1};
s_{i+1})$, and returns $(s_{i+1}, \comm_{i+1})$ to $\dev_{i+1}$, which
sends $\comm_{i+1}$ to $\dev_i$.

\emph{Transfer}\label{sec:stepS2S3S4}: $\dev_i$ notifies $\se_i$ that a
payment is under way; $\se_i$ locks token $\tau_i = (v, \comm_i)$. $\dev_i$
computes the audit ciphertext $\cipher_i \leftarrow \enc(\audpk, \uid_i;
\rho_i)$ and proof of correctness $\Gamma_i$, with $\instance = (\comm_i,
\cipher_i)$, $\wit = (\uid_i, s_i, \rho_i)$, and relation
\begin{equation*}
  \opn(\ck, \uid_i, s_i, \comm_i) = 1 \;\wedge\; \cipher_i = \enc(\audpk,
  \uid_i; \rho_i).
\end{equation*}
$\se_i$ and $\dev_i$ then jointly compute a signature of knowledge $\sigma_i$
for the relation
\begin{equation*}
  \sigverify(\rapk, \uid_i, \usk_i, \varphi_i) = 1 \;\wedge\; \opn(\ck,
  \uid_i, s_i, \comm_i) = 1,
\end{equation*}
where $\se_i$ contributes knowledge of $\usk_i$ and $\dev_i$ contributes the
remaining witness, following a sigma-protocol-based two-party protocol (see
Section~\ref{sec:sok-sol}). $\dev_i$ sends payment $p_i = (v, \comm_i,
\comm_{i+1}, \cipher_i, \sigma_i, \Gamma_i, \hist_i)$ to $\dev_{i+1}$.

\emph{Receipt}\label{sec:stepS5}: $\dev_{i+1}$ verifies $\hist_i$: if $i =
0$, it checks $\Sigma$, otherwise it verifies every $(\Sigma, \Gamma_k,
\sigma_k)_{k \in [i]}$. If $\Gamma_i$ and $\sigma_i$ also verify,
$\dev_{i+1}$ accepts and stores
\begin{equation*}
  T_{i+1} = (v, s_{i+1}, \comm_{i+1}, \hist_{i+1} = (\comm_i, \cipher_i,
  \sigma_i, \Gamma_i) \, \| \, \hist_i),
\end{equation*}
while $\se_{i+1}$ stores $\tau_{i+1} = (v, \comm_{i+1})$. Upon
acknowledgement, $\se_i$ deletes $\comm_i$.

\paragraph{Deposit and reconciliation.}\label{sec:deposit} User $\usr_n$
deposits token $\tau_n = (v, \comm_n)$ by instructing $\dev_n$, which informs
$\se_n$ to lock $\tau_n$; $\dev_n$ and $\se_n$ jointly compute a signature of
knowledge $\sigma_n$ over $\tau_n$, structured analogously to the payment
signatures above. $\dev_n$ sends deposit request $d_n = (\tau_n, \sigma_n,
\hist_n)$ to $\Int_n$, which forwards it to the central bank. The central
bank verifies $\sigma_n$ and $\hist_n$; if valid, it checks whether any
prior deposit shares the same origin $\comm_0$, and if so records double
spending and notifies the auditor, and otherwise credits $\Int_n$'s account,
which in turn credits $\usr_n$'s account. Upon notification, $\se_n$
deletes $\tau_n$. Appendix~\ref{sec:aggregate-deposits} extends the protocol
to support the aggregation of deposits.

\paragraph{Identifying double spenders.} For two conflicting deposit
requests $d_n$ and $d_m$ involving the same original token, let $\hist_n =
((\comm_k, \cipher_k, \sigma_k, \Gamma_k)_{k \in [n]}, \Sigma)$ and $\hist_m$
analogously. The auditor finds the divergence index $j$ such that $\comm_k =
\comm'_k$ for all $k \leq j$ but $\comm_{j+1} \neq \comm'_{j+1}$; the double
spender is the owner of $\tau_j = (v, \comm_j)$, identified by decrypting
$\cipher_j$ (or $\cipher'_j$).
